% GENERATED FILE. Edit canonical lessons, then run npm run generate:books.
% canonical-source-hash: fnv1a64:958b2f30204a7ad6
% canonical-lessons: PA-C117-kura, PA-C117-tokri, PA-C117-dabba, PA-C117-ghara, PA-C117-mombatti

\chapter{Rubbish, Basket, Box, Clay water pot, Candle}
\label{ch:pa-rubbish-basket-box-clay-water-pot-candle}
\hlchaptermodality{117}

\begin{chapteropening}
\textbf{By the end of this chapter:} \emph{I can say rubbish, a basket, a box, a clay water pot and a candle.}

Complete the last lesson of chapter 117: I can say rubbish, a basket, a box, a clay water pot and a candle.
\end{chapteropening}

\section[\texorpdfstring{\pa{ਕੂੜਾ} (kūṛā)}{ਕੂੜਾ (kūṛā)}]{\pa{ਕੂੜਾ} (kūṛā) — rubbish}

\label{lesson:PA-C117-kura}



\begin{quote}
Before the new one: say the Punjabi for a bucket, then the Punjabi for a broom.
\end{quote}

\subsection*{You'll want to know: \pa{ਕੂੜਾ}}
\textbf{\pa{ਕੂੜਾ}} — \emph{kūṛā} — \textquotedblleft{}rubbish\textquotedblright{}.

\begin{cousinweb}
\textbf{\pa{ਕੂੜਾ}} (\emph{kūṛā}) is a tadbhava, a word inherited from Sanskrit and reshaped by sound change. It comes from \emph{kūṭa}, \textquotedblleft{}heap, multitude\textquotedblright{} (CDIAL 3392). CDIAL takes the Sanskrit word, probably, from Dravidian; Punjabi \emph{kūṛā} is \textquotedblleft{}sweepings, dirt\textquotedblright{}.
\end{cousinweb}

\begin{grammarlens}[title={What you've built}]
One more word for this chapter.
\end{grammarlens}

\subsection*{Your turn}
\begin{itemize}
\raggedright
  \item \emph{Say it:} \emph{kūṛā}
  \item \emph{Say it:} \emph{kūṛā}, once more
  \item \emph{Say it:} say \emph{kūṛā}
  \item \emph{Recall:} say the Punjabi for a bucket, then the Punjabi for a broom, then say \emph{kūṛā} again
\end{itemize}

\begin{tcolorbox}[breakable,colback=teal!4,colframe=teal!35!black,title={Before you move on}]
What does \pa{ਕੂੜਾ} mean? (\textquotedblleft{}Rubbish\textquotedblright{}.) Say it once more.
\end{tcolorbox}
\section[\texorpdfstring{\pa{ਟੋਕਰੀ} (ṭokrī)}{ਟੋਕਰੀ (ṭokrī)}]{\pa{ਟੋਕਰੀ} (ṭokrī) — a basket}

\label{lesson:PA-C117-tokri}



\begin{quote}
Before the new one: say the Punjabi for a broom, then the Punjabi for rubbish.
\end{quote}

\subsection*{You'll want to know: \pa{ਟੋਕਰੀ}}
\textbf{\pa{ਟੋਕਰੀ}} — \emph{ṭokrī} — \textquotedblleft{}a basket\textquotedblright{}.

\begin{cousinweb}
\textbf{\pa{ਟੋਕਰੀ}} (\emph{ṭokrī}) is a tadbhava, a word inherited and reshaped by sound change. It comes from \emph{*ṭokka}, \textquotedblleft{}basket\textquotedblright{}, a form CDIAL reconstructs from its descendants (CDIAL 5477). CDIAL lists Punjabi \emph{ṭokrā}, with an added \emph{-r-}.
\end{cousinweb}

\begin{grammarlens}[title={What you've built}]
One more word for this chapter.
\end{grammarlens}

\subsection*{Your turn}
\begin{itemize}
\raggedright
  \item \emph{Say it:} \emph{ṭokrī}
  \item \emph{Say it:} \emph{ṭokrī}, once more
  \item \emph{Say it:} say \emph{ṭokrī}
  \item \emph{Recall:} say the Punjabi for a broom, then the Punjabi for rubbish, then say \emph{ṭokrī} again
\end{itemize}

\begin{tcolorbox}[breakable,colback=teal!4,colframe=teal!35!black,title={Before you move on}]
What does \pa{ਟੋਕਰੀ} mean? (\textquotedblleft{}A basket\textquotedblright{}.) Say it once more.
\end{tcolorbox}
\section[\texorpdfstring{\pa{ਡੱਬਾ} (ḍabbā)}{ਡੱਬਾ (ḍabbā)}]{\pa{ਡੱਬਾ} (ḍabbā) — a box}

\label{lesson:PA-C117-dabba}



\begin{quote}
Before the new one: say the Punjabi for rubbish, then the Punjabi for a basket.
\end{quote}

\subsection*{You'll want to know: \pa{ਡੱਬਾ}}
\textbf{\pa{ਡੱਬਾ}} — \emph{ḍabbā} — \textquotedblleft{}a box\textquotedblright{}.

\begin{cousinweb}
\textbf{\pa{ਡੱਬਾ}} (\emph{ḍabbā}) matches Hindi \emph{ḍabbā}, \textquotedblleft{}box\textquotedblright{}, which CDIAL lists under a reconstructed \emph{*ḍabba}, \textquotedblleft{}box\textquotedblright{}, a word it calls non-Aryan, perhaps Dravidian (CDIAL 5528). Its Punjabi forms there are \emph{ḍabb} and \emph{ḍibbā}.
\end{cousinweb}

\begin{grammarlens}[title={What you've built}]
One more word for this chapter.
\end{grammarlens}

\subsection*{Your turn}
\begin{itemize}
\raggedright
  \item \emph{Say it:} \emph{ḍabbā}
  \item \emph{Say it:} \emph{ḍabbā}, once more
  \item \emph{Say it:} say \emph{ḍabbā}
  \item \emph{Recall:} say the Punjabi for rubbish, then the Punjabi for a basket, then say \emph{ḍabbā} again
\end{itemize}

\begin{tcolorbox}[breakable,colback=teal!4,colframe=teal!35!black,title={Before you move on}]
What does \pa{ਡੱਬਾ} mean? (\textquotedblleft{}A box\textquotedblright{}.) Say it once more.
\end{tcolorbox}
\section[\texorpdfstring{\pa{ਘੜਾ} (ghaṛā)}{ਘੜਾ (ghaṛā)}]{\pa{ਘੜਾ} (ghaṛā) — a clay water pot}

\label{lesson:PA-C117-ghara}



\begin{quote}
Before the new one: say the Punjabi for a basket, then the Punjabi for a box.
\end{quote}

\subsection*{You'll want to know: \pa{ਘੜਾ}}
\textbf{\pa{ਘੜਾ}} — \emph{ghaṛā} — \textquotedblleft{}a clay water pot\textquotedblright{}.

\begin{cousinweb}
\textbf{\pa{ਘੜਾ}} (\emph{ghaṛā}) is a tadbhava, a word inherited from Sanskrit and reshaped by sound change. It comes from \emph{ghaṭa}, \textquotedblleft{}pot\textquotedblright{} (CDIAL 4406).
\end{cousinweb}

\begin{grammarlens}[title={What you've built}]
One more word for this chapter.
\end{grammarlens}

\subsection*{Your turn}
\begin{itemize}
\raggedright
  \item \emph{Say it:} \emph{ghaṛā}
  \item \emph{Say it:} \emph{ghaṛā}, once more
  \item \emph{Say it:} say \emph{ghaṛā}
  \item \emph{Recall:} say the Punjabi for a basket, then the Punjabi for a box, then say \emph{ghaṛā} again
\end{itemize}

\begin{tcolorbox}[breakable,colback=teal!4,colframe=teal!35!black,title={Before you move on}]
What does \pa{ਘੜਾ} mean? (\textquotedblleft{}A clay water pot\textquotedblright{}.) Say it once more.
\end{tcolorbox}
\section[\texorpdfstring{\pa{ਮੋਮਬੱਤੀ} (mombattī)}{ਮੋਮਬੱਤੀ (mombattī)}]{\pa{ਮੋਮਬੱਤੀ} (mombattī) — a candle}

\label{lesson:PA-C117-mombatti}



\begin{quote}
Before the new one: say the Punjabi for a box, then the Punjabi for a clay water pot.
\end{quote}

\subsection*{You'll want to know: \pa{ਮੋਮਬੱਤੀ}}
\textbf{\pa{ਮੋਮਬੱਤੀ}} — \emph{mombattī} — \textquotedblleft{}a candle\textquotedblright{}.

\begin{cousinweb}
\textbf{\pa{ਮੋਮਬੱਤੀ}} (\emph{mombattī}) joins words from two sources: \emph{mom}, \textquotedblleft{}wax\textquotedblright{}, usually traced to Persian \emph{mom}, and \emph{battī}, \textquotedblleft{}wick, lamp\textquotedblright{}, from Sanskrit \emph{varti} (CDIAL 11359). A candle is a \textquotedblleft{}wax lamp\textquotedblright{}.
\end{cousinweb}

\begin{grammarlens}[title={What you've built}]
One more word for this chapter.
\end{grammarlens}

\subsection*{Your turn}
\begin{itemize}
\raggedright
  \item \emph{Say it:} \emph{mombattī}
  \item \emph{Say it:} \emph{mombattī}, once more
  \item \emph{Say it:} say \emph{mombattī}
  \item \emph{Recall:} say the Punjabi for a box, then the Punjabi for a clay water pot, then say \emph{mombattī} again
\end{itemize}

\begin{tcolorbox}[breakable,colback=teal!4,colframe=teal!35!black,title={Before you move on}]
What does \pa{ਮੋਮਬੱਤੀ} mean? (\textquotedblleft{}A candle\textquotedblright{}.) Say it once more.
\end{tcolorbox}
